% ============================================================
% RESUMEN EJECUTIVO (Sección 0 — sin número de capítulo)
% ============================================================
\chapter*{Resumen Ejecutivo}
\addcontentsline{toc}{chapter}{Resumen Ejecutivo}
\markboth{Resumen Ejecutivo}{}

El Proyecto Demeter, desarrollado por la Asociaci\'on Universitaria ESIBot
en colaboraci\'on con la Universidad de Sevilla, ha alcanzado un estado
funcional documentado que demuestra la viabilidad t\'ecnica de una
plataforma IoT integral para investigaci\'on agr\'icola. A lo largo de
\textbf{veinte meses} desde su inicio en diciembre de 2024 --- de los
cuales solo cuarenta y cuatro d\'ias corresponden a desarrollo activo,
con 227 \textit{commits} concentrados entre enero y marzo de 2026 ---
el sistema ha llegado a cubrir las cuatro capas proyectadas
(firmware ESP32, \textit{edge computing} en Raspberry Pi, backend en
Google Cloud y frontend web) con documentaci\'on t\'ecnica completa de
84 p\'aginas y validaci\'on por capa mediante 140 tests unitarios y de
integraci\'on.

Sin embargo, el an\'alisis de sostenibilidad del proyecto en su forma
actual revela limitaciones estructurales que motivan su cierre formal:
(1) la complejidad multistack del sistema requiere un equipo
multidisciplinar sostenido en el tiempo que excede la capacidad actual
de la asociaci\'on; (2) los recursos econ\'omicos necesarios para el
mantenimiento operativo continuo superan las fuentes de financiaci\'on
disponibles; (3) el desarrollo se ha concentrado en un \'unico miembro,
configurando un punto de fallo estructural incompatible con la rotaci\'on
natural de una asociaci\'on estudiantil.

Este informe propone el cierre formal del proyecto en su forma actual
con preservaci\'on completa del conocimiento generado: publicaci\'on del
repositorio bajo licencia abierta, conservaci\'on de la documentaci\'on
t\'ecnica como caso de estudio, y disponibilidad para reactivaci\'on
parcial mediante TFG/TFM futuros sobre componentes espec\'ificos. Esta
estrategia permite cumplir los objetivos did\'acticos del proyecto,
preservar el valor acad\'emico generado, y mantener abierta la relaci\'on
de colaboraci\'on con la Universidad para futuras iniciativas.
